\chapter{Addresses, ports, routing, and DNS}
\chapterintro{A server is reachable only when several names and numbers agree: the hostname, IP address, transport port, listening socket, and route.}

\begin{objectives}
\begin{itemize}
  \item Explain the roles of IP addresses, transport ports, and socket endpoints.
  \item Distinguish loopback, private, public, unspecified, and wildcard addresses.
  \item Describe recursive DNS resolution and caching at a practical level.
  \item Explain routing and NAT without attributing application semantics to them.
\end{itemize}
\end{objectives}

\section{IP identifies an interface location}

An IP address identifies a network interface in a routing context. IPv4 writes
32 bits as four decimal octets, for example \code{192.0.2.10}. IPv6 uses 128
bits and hexadecimal groups, allowing compressed forms such as \code{2001:db8::10}.
The documentation ranges used here are reserved for examples and should not be
treated as real destinations.

An address is not a permanent identity for a person, application, or physical
machine. Interfaces can have several addresses; addresses can change; load
balancers can advertise one address for many servers; a mobile device can move
between networks. HTTP uses host and origin concepts above this layer precisely
because an IP address alone does not identify an application resource namespace.

\subsection{Address scopes that matter in the lab}

\begin{itemize}
  \item \code{127.0.0.1} is IPv4 loopback. Traffic stays inside the local host.
  \item \code{::1} is IPv6 loopback.
  \item Private IPv4 ranges are used inside networks and are not globally routed.
  \item Public addresses can be globally reachable subject to routing and policy.
  \item \code{0.0.0.0} used while binding means all local IPv4 interfaces; it is
        not a remote destination to which a client should normally connect.
\end{itemize}

Binding our server to loopback is a safety decision. A process listening on
\code{127.0.0.1:8080} is reachable by local clients but not directly through the
machine's other interfaces. Binding the same process to \code{0.0.0.0:8080}
may expose it to a local network, a container bridge, or more, depending on
firewall and routing rules.

\section{Ports identify transport endpoints}

An IP address gets traffic to a host interface. A 16-bit TCP or UDP port helps
the operating system deliver it to a transport endpoint. Port numbers range
from 0 through 65535. Ports below 1024 are commonly treated as privileged on
Unix-like systems. HTTP conventionally uses 80 and HTTPS 443, while development
servers often choose unprivileged ports such as 8000 or 8080.

A TCP connection is identified by protocol plus a four-tuple: local address,
local port, remote address, and remote port. This explains how a listening
server can support many clients on the same local port. Each accepted connection
has a different remote endpoint, and clients normally use temporary source ports.

\begin{mentalmodel}
The listening socket is a reception desk at one address and port. Accepting a
connection creates a separate conversation socket. The desk remains available
for new arrivals while workers handle existing conversations.
\end{mentalmodel}

\section{Routing chooses the next hop}

Hosts and routers maintain routing tables. A route associates an address prefix
with an outgoing interface and possibly a next-hop router. The most specific
matching prefix wins. If no more specific route exists, a default route usually
points toward a gateway.

Routers forward IP packets hop by hop. They do not establish one end-to-end
physical wire and do not need HTTP paths. Different packets can in principle
take different routes. IP makes a best effort; TCP above it recovers from loss
and reordering for applications that require a reliable stream.

\section{NAT rewrites addressing context}

Network Address Translation is common at boundaries between private and public
networks. A gateway can map many internal address-and-port pairs to one public
address with distinct ports. It records the mapping so reply traffic can be
directed back to the correct internal endpoint.

NAT is not a substitute for an application security policy, and it is not part
of HTTP. It changes the network view of endpoints. This can complicate inbound
connections, logging, and the interpretation of client addresses. Reverse
proxies add another layer: an application may see the proxy's TCP address while
the original address is described in a trusted HTTP field.

\section{DNS maps names to data}

The Domain Name System is a distributed, hierarchical database. A hostname can
have IPv4 address records (A), IPv6 address records (AAAA), aliases (CNAME), and
many other record types. DNS does not test whether an HTTP server is healthy; it
returns configured data and associated cache lifetimes.

A typical client asks a \term{recursive resolver}. The resolver may already
have a cached answer. Otherwise it follows delegations: root servers direct it
toward a top-level domain, those servers direct it toward authoritative servers,
and an authoritative server supplies the requested record. The resolver caches
answers according to their time-to-live values and returns them to the client.

\begin{figure}[htbp]
\centering
\begin{tikzpicture}[
  box/.style={draw=Rule,fill=PaperTint,rounded corners=1mm,
              minimum width=30mm,minimum height=10mm,font=\small\sffamily},
  arrow/.style={-{Latex[length=2mm]},Accent}
]
\node[box] (client) {client};
\node[box,right=10mm of client] (resolver) {recursive resolver};
\node[box,above right=8mm and 12mm of resolver] (root) {root};
\node[box,right=12mm of resolver] (tld) {TLD};
\node[box,below right=8mm and 12mm of resolver] (auth) {authoritative};
\draw[arrow] (client) -- node[above,font=\scriptsize]{question} (resolver);
\draw[arrow] (resolver) -- (root);
\draw[arrow] (resolver) -- (tld);
\draw[arrow] (resolver) -- (auth);
\draw[arrow,dashed] (resolver) -- node[below,font=\scriptsize]{cached answer} (client);
\end{tikzpicture}
\caption{A simplified recursive lookup. Referrals and replies are condensed.}
\label{fig:dns-resolution}
\end{figure}

\subsection{Why DNS answers vary}

DNS can return several addresses for redundancy or geographic distribution.
Different resolvers may have differently aged caches. A deployment can change
records, but old answers remain valid until their TTLs expire. A failed new
server therefore does not always affect every client at once, and a fixed DNS
record does not immediately repair clients holding an older answer.

\section{Names return at the HTTP layer}

After DNS, the hostname still matters. One IP address can host many HTTP origins.
In HTTP/1.1 the client sends a \code{Host} field so the server or reverse proxy
can select the intended origin. For HTTPS, the hostname is also used during TLS
identity verification and commonly during server-name negotiation. Replacing a
hostname with its address can therefore change both TLS and HTTP behavior.

\begin{pitfall}
``It pings'' does not prove an HTTP service works. Ping commonly uses ICMP, not
TCP port 8080, TLS, or HTTP. It can be blocked when HTTP works, or succeed when
no application is listening.
\end{pitfall}

\section{A disciplined diagnostic order}

When a server cannot be reached, ask questions from lower ambiguity to higher:

\begin{enumerate}
  \item What exact hostname, address, and port did the client choose?
  \item Did name resolution return the expected address family and address?
  \item Does the host have a route toward that destination?
  \item Is the server process listening on the intended local address and port?
  \item Do firewall or namespace boundaries permit the connection?
  \item Did TCP connect? If so, continue to TLS and HTTP evidence.
\end{enumerate}

\summarytitle

Addresses locate interfaces for IP forwarding; ports locate transport endpoints;
DNS maps names to configured data; routing chooses next hops; NAT translates
endpoint views. A server becomes reachable only when these mechanisms agree
with its bind address and the client's chosen destination.

\reviewtitle
\begin{enumerate}
  \item Why can many TCP clients share a server's port 8080?
  \item What is the practical difference between binding to loopback and to a wildcard address?
  \item Why might two clients temporarily receive different DNS answers for one name?
  \item Why can replacing a hostname with its IP address break HTTPS?
\end{enumerate}

\begin{exercise}[Endpoint inventory]
For a local development server, write the listening endpoint, one possible
client endpoint, and the resulting four-tuple. Repeat for two simultaneous
clients and underline what changes.
\end{exercise}

\begin{exercise}[Resolution narrative]
Narrate an uncached DNS lookup from a laptop to an authoritative answer. Then
repeat the story when every relevant answer is cached. Identify which network
round trips disappear.
\end{exercise}

